\documentclass[preview]{standalone}

\usepackage{amsmath}
\usepackage{amssymb}
\usepackage{stellar}
\usepackage{definitions}

\begin{document}

\id{psicologia-sonno-sogni-stati-alterati}
\genpage

\section{Sonno, sogni e stati alterati}

\begin{snippet}{sonno-introduzione}
    Il sonno è una profonda alterazione quotidiana della coscienza. Durante il
    sonno i muscoli si trovano in uno stato di quiescenza, mentre il cervello
    mantiene un'intensa attività.
\end{snippet}

\begin{snippetdefinition}{ritmo-circadiano-definition}{Ritmo circadiano}
    Il \emph{ritmo circadiano} è un ciclo temporale interno che regola livelli
    di attivazione, metabolismo, battito cardiaco e temperatura corporea,
    sincronizzandoli con l'orologio biologico dell'organismo.
\end{snippetdefinition}

\begin{snippet}{ritmi-circadiani-jet-lag}
    La luce solare contribuisce alla regolazione dei ritmi circadiani. Quando
    l'orologio interno si disallinea dal ciclo sonno-veglia, come nei turni
    notturni o nel jet lag, possono comparire difficoltà fisiche e cognitive.
\end{snippet}

\begin{snippetdefinition}{sonno-rem-definition}{Sonno REM}
    Il \emph{sonno REM} è una fase del sonno caratterizzata da movimenti oculari
    rapidi, aumento dell'attività elettrica cerebrale e sogni vividi.
\end{snippetdefinition}

\begin{snippet}{ciclo-sonno}
    Durante il sonno, l'attività cerebrale segue cicli regolari. Le fasi
    non-REM durano circa \(90\) minuti e includono il sonno profondo; la fase
    REM dura inizialmente circa \(10\) minuti ed è associata a sogni vividi. Un
    ciclo completo di circa \(100\) minuti si ripete dalle quattro alle sei
    volte per notte, con un aumento progressivo della durata delle fasi REM.
\end{snippet}

\includesnpt[width=62\%,src=/snippet/static-psychology/sleep-cycle-nrem-rem.jpg]{centered-img}

\begin{snippet}{funzioni-sonno}
    Il sonno non-REM è collegato a risparmio energetico e riposo. Il sonno REM
    sembra invece avere un ruolo nel consolidamento dei nuovi apprendimenti:
    quando una persona viene privata del sonno REM, nella notte successiva tende
    a recuperarne una quantità maggiore.
\end{snippet}

\begin{snippetdefinition}{deprivazione-sonno-definition}{Deprivazione di sonno}
    La \emph{deprivazione di sonno} è una condizione in cui una persona dorme
    troppo poco, con conseguenze negative sulle prestazioni cognitive e motorie.
\end{snippetdefinition}

\begin{snippetdefinition}{stati-coscienza-alterati-definition}{Stati di coscienza alterati}
    Gli \emph{stati di coscienza alterati} sono condizioni in cui l'esperienza
    consapevole differisce dallo stato ordinario di veglia, per esempio durante
    ipnosi, meditazione o pratiche rituali.
\end{snippetdefinition}

\begin{snippetdefinition}{ipnosi-definition}{Ipnosi}
    L'\emph{ipnosi} è uno stato di consapevolezza alternativo caratterizzato
    dalla capacità di rispondere a suggestioni con cambiamenti nella percezione,
    nella memoria, nella motivazione e nel senso di controllo di sé.
\end{snippetdefinition}

\begin{snippetdefinition}{induzione-ipnotica-definition}{Induzione ipnotica}
    L'\emph{induzione ipnotica} è l'insieme delle attività preliminari che
    riducono le distrazioni esterne e incoraggiano i partecipanti a concentrarsi
    sugli stimoli suggeriti, favorendo l'ingresso nello stato ipnotico.
\end{snippetdefinition}

\begin{snippetdefinition}{ipnotizzabilita-definition}{Ipnotizzabilità}
    L'\emph{ipnotizzabilità} è il grado di responsività di un individuo alla
    suggestione ipnotica.
\end{snippetdefinition}

\begin{snippetdefinition}{analgesia-ipnotica-definition}{Analgesia ipnotica}
    L'\emph{analgesia ipnotica} è il controllo del dolore ottenuto attraverso
    suggestioni ipnotiche, come immaginare la parte dolente separata dal resto
    del corpo o priva di sensibilità.
\end{snippetdefinition}

\begin{snippetdefinition}{meditazione-definition}{Meditazione}
    La \emph{meditazione} è una forma di alterazione della coscienza progettata
    per migliorare la conoscenza di sé e il benessere attraverso il
    raggiungimento di uno stato di profonda tranquillità.
\end{snippetdefinition}

\begin{snippet}{forme-meditazione}
    Nella \emph{meditazione concentrativa}, la persona regola il respiro,
    assume posture specifiche, riduce le stimolazioni esterne e concentra la
    mente. Nella \emph{meditazione mindfulness}, invece, impara a lasciare che
    ricordi e pensieri attraversino la mente senza reagire a essi.
\end{snippet}

\end{document}
